% =====================================================================
%  meta.tex — титул и вводная страница «как читать этот курс»
% =====================================================================

\thispagestyle{empty}

\begin{center}
  \vspace*{1.5cm}
  {\Huge\bfseries\color{cAccent} От электродинамики решётки\\[0.35em] до анализа остатков\par}
  \vspace{0.9cm}
  {\Large Учебное сопровождение проекта\par}
  \vspace{0.35em}
  {\Large\bfseries THz-Unified-Optimizer\par}
  \vspace{1.4cm}
  {\large\color{cInk2} Извлечение физических параметров проволочных поляризаторов\\
   из данных терагерцовой спектроскопии во временной области\par}
  \vspace{2.2cm}
  {\color{cInk2} Сессия D (статья и образовательный слой)\par}
  {\color{cInk2} \today\par}
\end{center}

\vfill

\begin{tcolorbox}[enhanced,colback=cPanel,colframe=cGrid,boxrule=0.8pt,arc=2pt]
\textbf{Кому адресован этот текст.} Заказчику и руководителю проекта, который ставит задачи и
принимает решения, но не является специалистом в поляризационной оптике, статистике нелинейных
подгонок и теории идентифицируемости. Поэтому курс \emph{обучает}, а не отчитывается: каждый термин
и символ определяются при первом употреблении, физическая интуиция даётся до формализма, а всякий
учебный тезис опирается на \emph{реальное} число проекта со ссылкой на породивший его файл.
Выдуманных учебных примеров здесь нет.
\end{tcolorbox}

\clearpage

\section*{Как читать этот курс}
\addcontentsline{toc}{section}{Как читать этот курс}
\markboth{Как читать этот курс}{}

\subsection*{Три части и почему именно в таком порядке}

\textbf{Часть~I — физика.} Что физически происходит с волной на решётке проводов и почему задача
сводится к двум комплексным числам на каждой частоте. Здесь же вводится различение, которое
определяет всю дальнейшую математику: терагерцовый прибор меряет \emph{поле}, а не интенсивность.

\textbf{Часть~II — оценивание.} Как из измерений извлекаются параметры. Начинаем с линейной
задачи, где всё решается замкнутой формулой и видно глазом, и лишь затем переходим к нелинейной,
где появляются множественные минимумы, мультистарт и вырождение. Этот порядок принципиален:
корреляция параметров вводится там, где её можно вычислить на бумаге.

\textbf{Часть~III — достоверность.} Чему верить в полученных числах. Правдоподобие,
информационные критерии, анализ остатков и главное различение всего курса: \emph{«модель лучше
описывает данные» и «параметр модели измерен» — разные утверждения.}

\subsection*{Сквозная идея}

\begin{tcolorbox}[enhanced,colback=cS1!7,colframe=cS1,boxrule=0pt,leftrule=3pt,arc=1pt]
Измерение содержит конечное число независимых чисел. Модель может содержать больше параметров, чем
измерение содержит чисел. Тогда параметры перестают существовать по отдельности~--- существуют
только их комбинации.
\end{tcolorbox}

\noindent
В этом проекте утверждение не метафорическое, а доказанное и с точной цифрой: угловая развёртка на
одной частоте способна ограничить \textbf{не более пяти вещественных чисел}. Отсюда прямо следует
почти всё содержание Частей~II и~III~--- в том числе то, почему один из ключевых параметров модели
приходится фиксировать, а не подгонять.

\subsection*{Условные обозначения врезок}

\begin{itemize}[leftmargin=1.4em,itemsep=0.3em]
  \item \textbf{\color{cAccent}Интуиция}~--- суть простыми словами, без формул. Если читать только
        эти врезки, получится связный рассказ без математики.
  \item \textbf{\color{cInk}Строго}~--- та же мысль в точной формулировке.
  \item \textbf{\color{cS3!70!black}Число проекта}~--- реальный результат с указанием артефакта,
        который его породил.
  \item \textbf{\color{cS2}Осторожно}~--- типичная ошибка или ограничение, о котором легко забыть.
\end{itemize}

\subsection*{Соглашение о мнимой единице}

Всюду используется инженерная мнимая единица~$\jj$ (а не~$i$), волна записывается как
$\propto e^{\jj(\omega t-kz)}$~--- так принято в радиофизике и так сделано в коде проекта. Это не
косметика: при обратном соглашении знаки всех фазовых задержек и мнимых частей импедансов меняются
на противоположные.
